% include file for chapter 2

\newpage \section{LITERÁRNÍ ČÁST} \label{lit}
Literární část této práce se bude zabývat výčtem současně užívaných technologií na výrobu kaprolaktamu (dále jen CPL), možnostmi zařízení na průmyslovou kapalinovou extrakci, matematickým popisem extraktoru a teorií rovnováhy kapalina-kapalina.

\subsection{Technologie kaprolaktamu} \label{CPLtechnology}
\textit{tohle je málo - rozvést alternativy (subsubsection)}

$\varepsilon$-Kaprolaktam je bílá krystalická látka s bodem tání 69,0°C~\cite{CPLSpolanaList}.
Tato sloučeina se sumárním vzorcem \ce{C6H11NO} je z chemického hlediska laktam $\varepsilon$-aminokapronové kyseliny, molekula tedy tvoří sedmičetný cyklus.
V literatuře je popsáno mnoho technologií na výrobu kaprolaktamu, z toho nejvýznamnější (?? \% světové produkce) je však postup popsaný v kapitole \ref{CHON}, který je zaveden i ve Spolaně a týká se jej tato práce.

\subsubsection{Tradiční výroba z cyklohexanonu} \label{CHON}
\textit{udělat reakční schéma, napsat proč je obsah síranu takový či makový}

Surovinou pro tuto technologie je cyklohexanon, typicky vyráběný oxidací cyklohexanu. Ten je nejprve substituován hydroxylaminem na cyklohexanon-oxim <<???CITACE: Ullmann???>> a následně podroben Beckmannově přesmyku katalyzovaném koncentrovanou kyselinou sírovou.
Vzniklá sloučenina kyseliny sírové a kaprolaktamu je nakonec neutralizována roztokem amoniaku, v důsledku čehož se reakční směs rozdělí na dvě fáze:
organickou, která obsahuje většinu kaprolaktamu, vedlejších produktů i část vody, a vodný roztok síranu amonného (dále jen SA) o koncentraci okolo 40~\%, který však vedle stopových nečistot obsahuje také zhruba 1~\% hm. CPL (v této práci bude složení uváděno jako hmotnostní zlomek, pokud nebude uvedeno jinak).

\subsection{Zařízení pro kapalinovou extrakci} \label{extraction}
\textit{stručný popis protiproudé extrakce, jaké jsou konstrukční možnosti, KARR UPSCALING}
zkratka RDC!

\subsection{Matematický popis extraktoru} \label{the_Math}
Kontinuální protiproudý extraktor v ustáleném stavu lze matematicky popsat buďto patrovým modelem nebo přenosovým modelem, tyto modely se však zásadně liší požadovanými informacemi i možností použití.
Patrový model, který bude dále popsán až v kapitole~\ref{compStage}, má využití zejména na porovnání separační účinnosti reálných extraktorů na základě změřených bilančních veličin, v případě Karrova extraktoru i na zvětšování měřítka z laboratorních experimentů na průmyslové zařízení, jak bylo popsáno v kapitole~\ref{extraction}.

Naproti tomu přenosový model k problému přistupuje jiným způsobem -- cílem je namodelovat chování extraktoru pouze ze znalostí základních vlastností, jedná se o vlastnosti látek, jako je hustota~($\rho$), dynamická viskozita~($\mu$) a difuzivita kaprolaktamu v dané látce~($D$), a také o vlastnosti celého systému, jako je mezifázové napětí a popis rovnováhy kapalina-kapalina v podobě rovnovážné funkce: $u=f(w)$.

Ve tomto textu budou dále používány m.j. tyto veličiny: $t$ znamená teplota v °C, $\dot{m}$ je hmotnostní tok, $Q$ je objemový tok, $u$ a $w$ jsou obecně hmotnostní zlomky CPL v organické fázi, resp. ve vodné fázi.
Veličiny s dolními indexy F, R, S a E se vztahují popořadě k surovině, rafinátu, solventu a extraktu.

\subsubsection{Přenosový model s axiální disperzí} \label{diff2}
V této kapitole bude podrobně vysvětlen výpočetní postup, který popisuje extraktor jako soustavu dvou nelineárních diferenciálních rovnic druhého řádu, kde každá rovnice popisuje jednu fázi.
Nezávislou veličinou modelu je bezrozměrná výška $z$ definovaná jako podíl výškové souřadnice $l$ vůči celkové výšce extraktoru $L$, a to ve směru od suroviny k rafinátu (či od extraktu k solventu).
Závislé veličiny jsou $c_x$ a $c_y$, tedy hmotnostní koncentrace CPL ve fázi ochuzované, resp. obohacované o CPL, což v případě této práce vždy znamená fázi vodnou, resp. organickou.
Tento výpočetní postup, který popsali S.~Hartland a A.~Kumar~\cite{HartlandKumar1997}, je vyjádřen soustavou rovnic~\ref{eqn:diff2}, kde $Pe_x$ a $Pe_y$ jsou Pécletova čísla (definováno v rovnici~\ref{eqn:Peclet} pro fázi $i$), $T_x$ je počet přenosových jednotek vzhledem k fázi x (definovaný v rovnici~\ref{eqn:Tx}) a $c_y^*$ je taková hodnota $c_x$, která by byla v rovnováze s $c_y$.
\begin{align}
\label{eqn:diff2}
	\frac{\D^2 c_x}{\D z^2} - Pe_x \frac{\D c_x}{\D z} &=  Pe_x T_x (c_x - c_y^*) \\
	\frac{\D^2 c_y}{\D z^2} + Pe_y \frac{\D c_y}{\D z} &= -Pe_y T_x (c_x - c_y^*) \frac{Q_x}{Q_y} \nonumber
\end{align}
\begin{equation}
\label{eqn:Peclet}
	Pe_i = \frac{Q_i L}{A E_i}
\end{equation}
$E_i$ v rovnici~\ref{eqn:Peclet} je koeficient axiální disperze pro fázi $i$.
$K_x$ v rovnici~\ref{eqn:Tx} je koeficient prostupu hmoty vztažený vždy k fázi $x$, $a$ je specifická mezifázová plocha a $A$ je plocha průřezu extraktoru.
\begin{equation}
	T_x = \frac{K_x a A L}{Q_x}
\label{eqn:Tx}
\end{equation}

Řešení soustavy rovnic~\ref{eqn:diff2} musí splňovat okrajové podmínky, které jsou pro bod $z=0$ definovány v rovnicích~\ref{eqn:bound2_1} a~\ref{eqn:bound2_2}.
\begin{align}
\label{eqn:bound2_1}
	\left( \frac{\D c_x}{\D z} \right)_{z=0} &= Pe_x (c_{x,0} - c_{x,\mathrm{in}}) \\
\label{eqn:bound2_2}
	\left( \frac{\D c_y}{\D z} \right)_{z=0} &= 0
\end{align}
V bodě $z=1$ musí být splněny podmínky definované v rovnicích~\ref{eqn:bound2_3} a~\ref{eqn:bound2_4}.
\begin{align}
\label{eqn:bound2_3}
	\left( \frac{\D c_x}{\D z} \right)_{z=1} &= 0 \\
\label{eqn:bound2_4}
	\left( \frac{\D c_y}{\D z} \right)_{z=1} &= Pe_y (c_{y,\mathrm{in}} - c_{y,1})
\end{align}
Řešením této soustavy pak je funkce $(c_x, c_y) = f(z)$, přičemž je nutné podotknout, že krajní hodnoty $c_x$ a $c_y$ odpovídají výstupním proudům, avšak nikoliv vstupním proudům, tedy $c_x(0) \neq c_\mathrm{F}$, $c_y(0) = c_\mathrm{E}$, $c_x(1) = c_\mathrm{R}$ a $c_y(1) \neq c_\mathrm{S}$. Při vstupech proudů dochází ke skokovým změnám koncentrací, lze tedy říci, že s vyšší hodnotou koeficientů axiální disperze se charakter toku obou fází vzdaluje pístovému toku a přibližuje ideálnímu mísiči.

Při výše uvedeném výpočtu je $L$ považováno za konstantu, takovýto výpočet tedy slouží jako výpočet simulační, zrovna tak jej však lze realizovat jako výpočet návrhový, kdy $L$ je v soustavě dif. rovnic neznámým parametrem, který je hledán na základě počátečního odhadu.
Při zavedení neznámého parametru je však nutné specifikovat další okrajovou podmínku, kterou je deklarace požadované koncentrace rafinátu (rovnice~\ref{eqn:boundExtra}).
\begin{equation}
	\left( c_x \right)_{z=1} = c_\mathrm{R}
\label{eqn:boundExtra}
\end{equation}

Programová implementace tohoto modelu je popsána v kapitole~\ref{compTransfer}.

\subsubsection{Přenosový model s pístovým tokem} \label{diff1}
Popsaný matematický model lze zjednodušit pomocí předpokladu pístového toku, pak jsou koeficienty axiální disperze nulové a lze odvodit soustavu rovnic~\ref{eqn:diff1}.
\begin{align}
\label{eqn:diff1}
	\frac{\D c_x}{\D z} &= -T_x (c_x - c_y^*)\\
	\frac{\D c_y}{\D z} &= -T_x (c_x - c_y^*) \frac{Q_x}{Q_y} \nonumber
\end{align}
Při pístovém toku jsou však odlišné okrajové podmínky -- nedochází ke skokovým změnám koncentrací, okrajové podmínky jsou tedy vlastně deklaracemi krajních hodnot $c_x$ a $c_y$, jak je uvedeno v rovnicích~\ref{eqn:bound1_1}, ~\ref{eqn:bound1_2}.
Při návrhovém výpočtu je zde, obdobně jako u modelu s axiální disperzí, jedna okrajová podmínka navíc, která již byla definována v rovnici~\ref{eqn:boundExtra}.
\begin{align}
\label{eqn:bound1_1}
	\left( c_x \right)_{z=0} &= c_\mathrm{F} \\
\label{eqn:bound1_2}
	\left( c_y \right)_{z=1} &= c_\mathrm{S}
\end{align}

Model s pístovým tokem je využit jako mezivýpočet při implementaci axiálně disperzního modelu, jak je popsáno v kapitole~\ref{compTransfer}.

\subsection{Charakterizace přenosu hmoty} \label{insfun}
Výpočet přenosového modelu, jak byl popsán v kapitole~\ref{the_Math}, vyžaduje znalost průřezu extraktoru $A$ a veličin charakterizujících přenos hmoty v extraktoru, a to $K_x$, $E_y$, $E_x$ a $a$.
Pro přílišnou složitost hydrodynamických výpočtů nelze tyto veličiny určit jinak než na základě empirických rovnic.
Ty jsou v literatuře dostupné pro různé konstrukce extraktorů, přičemž v této práci byly převzaty rovnice popisující výplňovou kolonu, RDC, Kühniho a Karrův extraktor.
Tyto empirické vztahy pracují se základními vlastnostmi obou fází ($\rho$, $\mu$ a $D$), mezifázovým napětím $\gamma$ a parametry specifickými pro konstrukci extraktoru, což pro výplňovou kolonu je popis výplně (tvar výplně, charakteristický rozměr $d_\mathrm{P}$, specifická plocha $a_\mathrm{P}$ a porozita $\varepsilon$), pro RDC a Kühniho extraktor je to pouze frekvence $N$ a pro Karrův extraktor jsou parametry hlavně frekvence $N$ a zdvih $s$, případně i relativní volná plocha disku $\sigma$ a vzdálenost mezi disky $H$.

Výpočet koeficientu prostupu hmoty $K_x$ není přímý, skládá se totiž z koeficientů přestupu obou fází ($\beta$), jak je popsáno v rovnici~\ref{eqn:Kx}~\cite{HartlandKumar1997}, kde $m$ je směrnice rovnovážné funkce $c_y = f(c_x)$ na dané souřadnici kolony.
V některých případech je k dispozici výpočet koeficientů axiální disperze, jindy je namísto toho vypočteno přímo Pécletovo číslo (viz rovnici~\ref{eqn:Peclet}).
Specifickou mezifázovou plochu $a$ lze vypočítat ze zádrže dispergované fáze $\phi$ a Sauterovy střední velikosti kapek $d_{32}$, jak je popsáno v rovnici~\ref{eqn:aint}~\cite{HartlandKumar1997}.
Porozita je přitom parametrem pouze pro výplňovou kolonu, zatímco pro ostatní konstrukce extraktoru se předpokládá jednotková porozita, míchadla i hřídel se do ní tedy nezapočítávají.
Rovnice~\ref{eqn:R_ratio}"-\ref{eqn:Sherwood}~\cite{HartlandKumar1997} pak definují různé mezivýsledky, které budou v následujících podkapitolách často používány.
Popořadě se jedná o poměr toků dispergované a kontinuální fáze $R$, rychlost kluzu fází (z angl. slip velocity) $v_\mathrm{s}$, střední hustotu $\bar{\rho}$ i rozdíl hustot $\Delta\rho$, Reynoldsovo číslo pro kontinuální fázi (charakterizující její turbulenci), Schmidtovo číslo pro fázi $i$ (charakterizující difúzi a vnitřní tření v kapalině) a konečně Sherwoodovo číslo pro fázi $i$, z nějž lze určit hodnoty $\beta_i$.

\begin{align}
\label{eqn:Kx}
	\frac{1}{K_x} &= \frac{1}{\beta_x} + \frac{1}{ \beta_y m } \\
\label{eqn:aint}
	a &= \frac{6 \varepsilon \phi }{ d_{32} } \\
\label{eqn:R_ratio}
	R &= \frac{ v_\mathrm{D} }{ v_\mathrm{C} } \\
	v_\mathrm{s} &= \frac{ v_\mathrm{C} }{ \varepsilon (1-\phi) } + \frac{ v_\mathrm{D} }{ \varepsilon \phi} \\
	\bar{\rho} &= \rho_\mathrm{C} (1-\phi) + \rho_\mathrm{D} \phi \\
	\Delta\rho &= |\rho_\mathrm{C} - \rho_\mathrm{D}| \\
	Re_\mathrm{C} &= \frac{ \rho_\mathrm{C} d_{32} v_\mathrm{s} }{ \mu_\mathrm{C} } \\
	Sc_i &= \frac{ \mu_i }{ \rho_i D_i } \\
\label{eqn:Sherwood}
	Sh_i &= \frac{ \beta_i d_{32} }{ D_i } \\
\label{eqn:chi}
	\chi &= \frac{\phi}{1 - \phi}  \left(\frac{\rho_\mathrm{D} \mu_\mathrm{C} }{ \rho_\mathrm{C} \mu_\mathrm{D}}\right)^{0.3}
\end{align}
Užitečnou dodatkovou informaci můžeme nalézt v rovnici~\ref{eqn:chi}, kde je definován parametr $\chi$, který předpovídá stabilitu disperze~\cite{chi} (strana~149).
Hodnota menší než 0,3 zaručuje stabilitu dispergované fáze, hodnota do 0,5 stále znamená nejspíše stabilní disperzi, oblast 0,5--2,0 je oblast neurčitá, hodnota větší než 2,0 pak předpovídá možnou inverzi fází, hodnota větší než 3,3 ji zaručuje.

Přenosový model vyžaduje také znalost průřezu extraktoru -- ten je při simulačním výpočtu zadán, zatímco při návrhovém výpočtu je určen pomocí meze zahlcení $v_\mathrm{C,E}$, což je charakteristická rychlost, při jejímž překročení dochází ke značnému únosu dispergované fáze fází kontinuální~\cite{Perry_15} (strana~15"-11).
Je tedy určen takový průřez, aby rychlost kontinuální fáze $v_\mathrm{C}$ dosahovala hodnoty meze zahlcení redukované na určitý bezpečnostní faktor, obvykle na 50--70~\%.

\subsubsection{Výpočty pro výplňovou kolonu} \label{insfunPacked}

\subsubsection{Výpočty pro RDC a Kühniho extraktor} \label{insfunRotDisc}

\subsubsection{Výpočty pro Karrův extraktor} \label{insfunKarr}


\subsection{Rovnováha kapalina-kapalina} \label{LLEtheory}
\textit{Krátký teoretický popis, proč my na teorii kašleme, nějaká data z literatury, ta ale nestačí a proč}
